% Faithful editable transcription of official CCC 2024 article 1, pp. 1:8--1:13.
% Typography only: author numbering is carried in unnumbered headings.
\section*{Definitions and notation (author Section 2)}
We write $\mathbb N$ to denote the set of non-negative integers.
In the introduction, we worked with functions $f:\{0,1\}^n\to\{0,1\}$. Going forward, it will be more convenient to encode a bit $b\in\{0,1\}$ as the value $(-1)^b$. Thus, we will work with functions $f:\{\pm1\}^n\to\{\pm1\}$. We continue to use the notation $f^{\oplus t}$, but now $f^{\oplus t}$ denotes the product of $t$ copies of $f$ on independent inputs.
We use the following notation to describe decision trees.
\paragraph{Definition 2.1 (Decision trees).} For a function $f:\{\pm1\}^n\to\{\pm1\}$, we define $\DT(f)$ to be the minimum depth of a decision tree computing $f$. In the other direction, for a parameter $D\in\mathbb N$, we define $\DT[D]$ to be the class of all functions $f:\{\pm1\}^n\to\{\pm1\}$ that can be computed by depth-$D$ decision trees. (The parameter $n$ will always be clear from context.)
We denote random variables using boldface. We write $\mathbf x\sim\mu$ to indicate that the random variable $\mathbf x$ is sampled from the distribution $\mu$.
\paragraph{Definition 2.4 (Tensor product of probability distributions).} Let $\mu_1,\ldots,\mu_t$ be probability distributions over the spaces $\Omega_1,\ldots,\Omega_t$. Sample $\mathbf x^1\sim\mu_1,\ldots,\mathbf x^t\sim\mu_t$ independently. The tensor product $\mu_1\otimes\cdots\otimes\mu_t$ is the probability distribution of $(\mathbf x^1,\ldots,\mathbf x^t)$. As a special case, we define
\[\mu^{\otimes t}=\underbrace{\mu\otimes\mu\otimes\cdots\otimes\mu}_{t\text{ copies}}.\]
We use the following standard definition to reason about average-case hardness of $\{\pm1\}$-valued functions.
\paragraph{Definition 2.5 (Correlation).} Let $g,h:\{\pm1\}^n\to\mathbb R$ be functions and let $\mu$ be a distribution over $\{\pm1\}^n$. We define
\[\Corr_\mu(g,h)=\mathbb E_{\mathbf x\sim\mu}[g(\mathbf x)\cdot h(\mathbf x)].\]
More generally, if $\mathcal C$ is a class of functions $g:\{\pm1\}^n\to\mathbb R$, then we define
\[\Corr_\mu(\mathcal C,h)=\max_{g\in\mathcal C}\Corr_\mu(g,h).\]
If $\mu$ is omitted, then by default it is assumed to be the uniform distribution over $\{\pm1\}^n$.
If $g$ and $h$ are $\{\pm1\}$-valued, then a bound $|\Corr(g,h)|\le\varepsilon$ is equivalent to the statement that $g$ agrees with $h$ on at most a $(1/2+\varepsilon/2)$-fraction of inputs, because for any two $\{0,1\}$-valued random variables $\mathbf a,\mathbf b$, we have $\Pr[\mathbf a=\mathbf b]=\frac12+\frac12\mathbb E[(-1)^{\mathbf a}\cdot(-1)^{\mathbf b}]$.
\paragraph{Definition 2.9 (Log-concave).} A function $f:[0,\infty)\to(0,\infty)$ is log-concave if $\log f$ is concave, i.e., for every $x,y\in[0,\infty)$ and $\lambda\in(0,1)$, we have $f(x)^\lambda\cdot f(y)^{1-\lambda}\le f(\lambda x+(1-\lambda)y)$.
If $f$ is log-concave, then by induction on $t$, we have $\prod_{i=1}^t f(x_i)\le f(\bar x)^t$ where $\bar x=\frac1t\sum_{i=1}^t x_i$.
% STATEMENT_BEGIN
\paragraph{Lemma 3.2 (XOR lemma for decision trees under robust hardness assumptions, general version).}
Let $h_1,\ldots,h_t:\{\pm1\}^r\to\{\pm1\}$ be functions, and define $h(y^{(1)},\ldots,y^{(t)})=\prod_{i=1}^t h_i(y^{(i)})$. Let $\mu_1,\ldots,\mu_t$ be distributions over $\{\pm1\}^r$, and define $\mu=\mu_1\otimes\cdots\otimes\mu_t$. Let $\varepsilon:[0,\infty)\to(0,\infty)$ be a log-concave function, and assume that for every $i\in[t]$ and every $D\in\mathbb N$, we have
\[\Corr_{\mu_i}(h_i,\DT[D])\le\varepsilon(D).\]
Then for every $D\in\mathbb N$, we have
\[\Corr_\mu(h,\DT[Dt/2])\le O(\varepsilon(D))^t.\]
% STATEMENT_END
% PROOF_BEGIN
The first step of the proof of Lemma 3.2 is the following claim, which enables us to relate the success probability of a tree to the success probabilities of its subtrees.
\paragraph{Claim 3.3 (Law of total correlation).} Let $h,T,E:\{\pm1\}^r\to\{\pm1\}$. Let $\mu$ be a distribution over $\{0,1\}^r$. For each $b\in\{\pm1\}$, let $p_b=\Pr_{\mathbf y\sim\mu}[E(\mathbf y)=b]$, and let $\mu_b$ be the conditional distribution $(\mathbf y\sim\mu\mid E(\mathbf y)=b)$. Suppose that $T$ can be decomposed in the form
\[T(y)=\begin{cases}T_{+1}(y)&\text{if }E(y)=+1,\\T_{-1}(y)&\text{if }E(y)=-1\end{cases}\]
for some $T_{+1},T_{-1}:\{\pm1\}^r\to\{\pm1\}$. Then
\[\Corr_\mu(h,T)=\sum_{b\in\{\pm1\}}p_b\cdot\Corr_{\mu_b}(h,T_b).\]
\begin{proof}
\begin{align*}
\Corr_\mu(h,T)&=\mathbb E_{\mathbf y\sim\mu}[h(\mathbf y)\cdot T(\mathbf y)]\\
&=\sum_{b\in\{\pm1\}}p_b\cdot\mathbb E_{\mathbf y\sim\mu}[h(\mathbf y)\cdot T(\mathbf y)\mid E(\mathbf y)=b] &&\text{(Law of total expectation)}\\
&=\sum_{b\in\{\pm1\}}p_b\mathbb E_{\mathbf y\sim\mu_b}[h(\mathbf y)\cdot T_b(\mathbf y)].
\end{align*}
\end{proof}
Next, we consider the following notion of ``fair'' decision trees due to Shaltiel [61].
\paragraph{Definition 3.4 ($(D_1,\ldots,D_t)$-fair decision trees [61]).} Let $T:\{\pm1\}^{rt}\to\{\pm1\}$ be a decision tree and let $D_1,\ldots,D_t\in\mathbb N$. We say that $T$ is $(D_1,\ldots,D_t)$-fair if for every input $\vec y=(y^{(1)},\ldots,y^{(t)})\in(\{\pm1\}^r)^t$, for every $i\in[t]$, the computation $T(\vec y)$ makes at most $D_i$ queries to $y^{(i)}$.
The key to proving Lemma 3.2 is to generalize Definition 3.4 to the case of a set of tuples $(D_1,\ldots,D_t)$.
\paragraph{Definition 3.5 ($Q$-fair decision trees).} Let $T:\{\pm1\}^{rt}\to\{\pm1\}$ be a decision tree and let $Q\subseteq\mathbb N^t$. We say that $T$ is $Q$-fair if for every input $\vec y=(y^{(1)},\ldots,y^{(t)})\in(\{\pm1\}^r)^t$, there is some tuple $(D_1,\ldots,D_t)\in Q$ such that for every $i\in[t]$, the computation $T(\vec y)$ makes at most $D_i$ queries to $y^{(i)}$.
We emphasize that the tuple $(D_1,\ldots,D_t)$ is permitted to vary from one input $\vec y$ to another. Therefore, the fact that a tree is $Q$-fair does not necessarily imply that there is some $(D_1,\ldots,D_t)\in Q$ such that the tree is $(D_1,\ldots,D_t)$-fair. Given the concept of $Q$-fairness, it is relatively straightforward to prove the following claim by induction on the depth of $T$. The claim generalizes the analysis by Shaltiel [61], who considered the case of $(D_1,\ldots,D_t)$-fair decision trees and focused on the uniform distribution.
\paragraph{Claim 3.6 (XOR lemma for $Q$-fair decision trees).} Let $h_1,\ldots,h_t:\{\pm1\}^r\to\{\pm1\}$ be functions, and define $h(y^{(1)},\ldots,y^{(t)})=\prod_{i=1}^t h_i(y^{(i)})$. Let $\mu_1,\ldots,\mu_t$ be distributions over $\{\pm1\}^r$, and define $\mu=\mu_1\otimes\cdots\otimes\mu_t$. Let $Q\subseteq\mathbb N^t$ and let $T:\{\pm1\}^{rt}\to\{\pm1\}$ be a $Q$-fair decision tree. Then
\[\Corr_\mu(h,T)\le\sum_{(D_1,\ldots,D_t)\in Q}\prod_{i=1}^t\Corr_{\mu_i}(h_i,\DT[D_i]).\]
\begin{proof}
Assume without loss of generality that $T$ never queries the same variable twice. For the base case, if $T$ has depth 0, then $T$ is a constant function, so
\[|\Corr_\mu(h,T)|=\prod_{i=1}^t\left|\mathbb E_{\mathbf y^{(i)}\sim\mu_i}[h_i(\mathbf y^{(i)})]\right|=\prod_{i=1}^t\Corr_{\mu_i}(h_i,\DT[0]).\]
Since $T$ is $Q$-fair, $Q$ must be nonempty. The lemma follows because $\Corr_{\mu_i}(h_i,\DT[0])\le\Corr_{\mu_i}(h_i,\DT[D_i])$ for every $D_i\in\mathbb N$. For the inductive step, let $y_{j^*}^{(i^*)}$ be the variable queried by the root of the tree. Let $T_{+1}$ and $T_{-1}$ be the children of the root, corresponding to the cases $y_{j^*}^{(i^*)}=+1$ and $y_{j^*}^{(i^*)}=-1$ respectively. Define
\[Q'=\{(D_1,\ldots,D_{i^*-1},D_{i^*}-1,D_{i^*+1},\ldots,D_t):(D_1,\ldots,D_t)\in Q\text{ and }D_{i^*}\ne0\}.\]
Then $T_{+1}$ and $T_{-1}$ are both $Q'$-fair.
For each $b\in\{\pm1\}$, define
\[p_b=\Pr_{\mathbf y^{(i^*)}\sim\mu_{i^*}}[\mathbf y_{j^*}^{(i^*)}=b].\]
Let $\mu_{i^*}^b$ be the conditional distribution $(\mathbf y^{(i^*)}\sim\mu_{i^*}\mid\mathbf y_{j^*}^{(i^*)}=b)$, and for $i\ne i^*$, let $\mu_i^b=\mu_i$. Let $\mu^b=\mu_1^b\otimes\cdots\otimes\mu_t^b$. By Claim 3.3 and the induction hypothesis, we have
\begin{align*}
\Corr_\mu(h,T)&=\sum_{b\in\{\pm1\}}p_b\cdot\Corr_{\mu^b}(h,T_b)\\
&\le\sum_{b\in\{\pm1\}}p_b\cdot\sum_{(D_1,\ldots,D_t)\in Q'}\prod_{i=1}^t\Corr_{\mu_i^b}(h_i,\DT[D_i])\\
&=\sum_{(D_1,\ldots,D_t)\in Q'}\left(\sum_{b\in\{\pm1\}}p_b\cdot\Corr_{\mu_{i^*}^b}(h_{i^*},\DT[D_{i^*}])\right)\cdot\Pi_{\ne i^*}(D_1,\ldots,D_t)
\end{align*}
where $\Pi_{\ne i^*}(D_1,\ldots,D_t)=\prod_{i\in[t],i\ne i^*}\Corr_{\mu_i}(h_i,\DT[D_i])$. Now we bound the inner sum. By Claim 3.3, for any $D_{i^*}$, we have
\[\Corr_{\mu_{i^*}}(h_{i^*},\DT[D_{i^*}+1])\ge\sum_{b\in\{\pm1\}}p_b\cdot\Corr_{\mu_{i^*}^b}(h_{i^*},\DT[D_{i^*}]),\]
because we can approximate $h_{i^*}$ with respect to $\mu_{i^*}$ by first querying $y_{j^*}^{(i^*)}$ and then using optimal subtrees of depth $D_{i^*}$. For every $(D_1,\ldots,D_t)\in Q'$, we have $(D_1,\ldots,D_{i^*-1},D_{i^*}+1,D_{i^*+1},\ldots,D_t)\in Q$. Therefore,
\[\Corr_\mu(h,T)\le\sum_{(D_1,\ldots,D_t)\in Q}\prod_{i=1}^t\Corr_{\mu_i}(h_i,\DT[D_i]).\]
\end{proof}
Given Claim 3.6, our XOR lemma for decision trees under a robust hardness assumption (Lemma 3.2) readily follows, as we now show.
\begin{proof}[Proof of Lemma 3.2]
Let $T:\{\pm1\}^{rt}\to\{\pm1\}$ be a decision tree of depth at most $Dt/2$. Let $Q$ be the set of $t$-tuples $(D_1,\ldots,D_t)\in\mathbb N^t$ such that (1) $D_1+\cdots+D_t\le Dt$ and (2) $D_i$ is an integer multiple of $\lceil D/2\rceil$ for every $i$. We claim that $T$ is $Q$-fair. Indeed, let $\vec y=(y^{(1)},\ldots,y^{(t)})$ be any input, and let $D_i$ be the number of queries that $T(\vec y)$ makes to $y^{(i)}$. Let $D_i'$ be the smallest integer multiple of $\lceil D/2\rceil$ such that $D_i\le D_i'$. Then $D_i'\le D_i+(\lceil D/2\rceil-1)$, and hence $D_1'+\cdots+D_t'\le Dt/2+t\cdot(\lceil D/2\rceil-1)\le Dt$, showing that $(D_1',\ldots,D_t')\in Q$.
Therefore, by Claim 3.6,
\[\Corr_\mu(h,T)\le\sum_{(D_1,\ldots,D_t)\in Q}\prod_{i=1}^t\Corr_{\mu_i}(h_i,\DT[D_i]).\]
For any $(D_1,\ldots,D_t)\in Q$, we can define $(D_1',\ldots,D_t')$ such that $D_i'\ge D_i$ and $D_1'+\cdots+D_t'$ is exactly $Dt$ rather than being at most $Dt$. Then $\Corr_{\mu_i}(h_i,\DT[D_i])\le\Corr_{\mu_i}(h_i,\DT[D_i'])$, so
\begin{align*}
\Corr_\mu(h,T)&\le\sum_{(D_1,\ldots,D_t)\in Q}\prod_{i=1}^t\Corr_{\mu_i}(h_i,\DT[D_i'])\\
&\le\sum_{(D_1,\ldots,D_t)\in Q}\prod_{i=1}^t\varepsilon(D_i')\\
&\le\sum_{(D_1,\ldots,D_t)\in Q}\varepsilon(D)^t&&\text{(Log-concavity)}\\
&=|Q|\cdot\varepsilon(D)^t.
\end{align*}
To bound $|Q|$, observe that if $(D_1,\ldots,D_t)\in Q$, then we can write $D_i=c_i\cdot\lceil D/2\rceil$ for some nonnegative integers $c_1,\ldots,c_t$. Furthermore, $Dt\ge\sum_i c_i\cdot\lceil D/2\rceil\ge(D/2)\cdot\sum_i c_i$, so $c_1+\cdots+c_t\le2t$. Therefore, $|Q|$ is at most the number of ways that $2t$ can be partitioned into $t+1$ nonnegative integers, which is precisely $\binom{3t}{t}$. Thus,
\[\Corr_\mu(h,T)\le\binom{3t}{t}\cdot\varepsilon(D)^t\le O(\varepsilon(D))^t.\]
\end{proof}
% PROOF_END
